% Part: first-order-logic
% Chapter: natural-deduction
% Section: proving-things-quant

\documentclass[../../../include/open-logic-section]{subfiles}

\begin{document}

\olfileid{fol}{ntd}{prq}

\olsection{带量词的推导}

\begin{ex}
处理量词时，我们必须确保不违反本征变元条件，而这有时需要我们调整某些推理的进行顺序。一般来说，先处理那些受本征变元条件约束的规则会很有帮助（它们会出现在完成后的证明中靠近底部的位置）。

下面我们来看看如何给出公式 $\lexists[x][\lnot !A(x)] \lif \lnot \lforall[x][!A(x)]$ 的一个推导。
像往常一样，我们这样开始：
\begin{prooftree}
\AxiomC{}
\UnaryInfC{$\lexists[x][\lnot !A(x)]\lif \lnot \lforall[x][!A(x)]$}
\end{prooftree}
我们首先写下，若要使用 \Intro{\lif} 规则得出最后一步，需要什么前提。
\begin{prooftree}
\AxiomC{$\Discharge{\lexists[x][\lnot !A(x)]}{1}$}
\DeduceC{$\lnot \lforall[x][!A(x)]$}
\DischargeRule{\Intro{\lif}}{1}
\UnaryInfC{$\lexists[x][\lnot !A(x)]\lif \lnot \lforall[x][!A(x)]$}
\end{prooftree}
由于对 $\lnot \lforall[x][!A(x)]$ 没有明显的可用规则，我们将继续构造推导，以便使用 \Elim{\lexists} 规则。此处我们必须关注本征变元条件，并选择一个不在 $\lexists[x][!A(x)]$ 中、也不在其所依赖的任何假设中出现的常项符号。
（不过，既然目前没有出现任何常项符号，那么怎么选择都可以。）
\begin{prooftree}
\AxiomC{$\Discharge{\lexists[x][\lnot !A(x)]}{1}$}
\AxiomC{$\Discharge{\lnot !A(a)}{2}$}
\DeduceC{$\lnot \lforall[x][!A(x)]$}
\DischargeRule{\Elim{\lexists}}{2}
\BinaryInfC{$\lnot \lforall[x][!A(x)]$}
\DischargeRule{\Intro{\lif}}{1}
\UnaryInfC{$\lexists[x][\lnot !A(x)] \lif \lnot \lforall[x][!A(x)]$}
\end{prooftree}
为了推导出 $\lnot \lforall[x][!A(x)]$，我们将尝试使用 \Intro{\lnot} 规则：这要求我们推导出一个矛盾，可能需要使用 $\lforall[x][!A(x)]$ 作为额外假设。当然，这个矛盾可能涉及假设 $\lnot !A(a)$，而该假设将被 \Elim{\lexists} 推理解除。我们可以如下设置：
\begin{prooftree}
\AxiomC{$\Discharge{\lexists[x][\lnot !A(x)]}{1}$}
\AxiomC{$\Discharge{\lnot !A(a)}{2}, \Discharge{\lforall[x][!A(x)]}{3}$}
\DeduceC{$\lfalse$}
\DischargeRule{\Intro{\lnot}}{3}
\UnaryInfC{$\lnot \lforall[x][!A(x)]$}
\DischargeRule{\Elim{\lexists}}{2}
\BinaryInfC{$\lnot \lforall[x][!A(x)]$}
\DischargeRule{\Intro{\lif}}{1}
\UnaryInfC{$\lexists[x][\lnot !A(x)]\lif \lnot \lforall[x][!A(x)]$}
\end{prooftree}
看起来我们已经接近得出矛盾了。最容易应用的规则是 \Elim{\lforall}，它没有本征变元条件。由于我们可以使用任何想要的项来替换全称量化的~$x$，因此最合理的做法是继续使用~$a$，这样我们就能得出矛盾。
\begin{prooftree}
\AxiomC{$\Discharge{\lexists[x][\lnot !A(x)]}{1}$}
\AxiomC{$\Discharge{\lnot !A(a)}{2}$}
\AxiomC{$\Discharge{\lforall[x][!A(x)]}{3}$}
\RightLabel{\Elim{\lforall}}
\UnaryInfC{$!A(a)$}
\RightLabel{\Elim{\lnot}}
\BinaryInfC{$\lfalse$}
\DischargeRule{\Intro{\lnot}}{3}
\UnaryInfC{$\lnot \lforall[x][!A(x)]$}
\DischargeRule{\Elim{\lexists}}{2}
\BinaryInfC{$\lnot \lforall[x][!A(x)]$}
\DischargeRule{\Intro{\lif}}{1}
\UnaryInfC{$\lexists[x][\lnot !A(x)]\lif \lnot \lforall[x][!A(x)]$}
\end{prooftree}

在此时复查本征变元条件是否被违反尤为重要，特别是在处理量词时。由于我们应用的唯一受制于本征变元条件的规则是 \Elim{\exists}，且本征变元~$a$ 并未出现在它所依赖的任何假设中，因此这是一个正确的推导。
\end{ex}

\begin{ex}
有时我们需要从另一些公式推导出一个公式。在这些情况下，我们可能会有未解除的假设。记录好我们的假设以及最终目标是非常重要的。

我们来看看如何从假设 $\lexists[x][(!A(x)
\land !B(x))]$ 和 $\lforall[x][(!B(x) \lif !C(x,b))]$ 推导出公式 $\lexists[x][!C(x,b)]$。像往常一样，我们将结论写在底部。
\begin{prooftree}
\AxiomC{}
\UnaryInfC{$\lexists[x][!C(x,b)]$}
\end{prooftree}

我们有两个前提可以利用。要使用第一个前提，即试图从 $\lexists[x][(!A(x)
  \land !B(x))]$ 推导出 $\lexists[x][!C(x, b)]$，我们会用到 \Elim{\lexists} 规则。由于它受本征变元条件约束，我们将首先应用该规则。我们得到如下推导：
\begin{prooftree}
\AxiomC{$\lexists[x][(!A(x) \land !B(x))]$}
\AxiomC{$\Discharge{!A(a) \land !B(a)}{1}$}
\DeduceC{$\lexists[x][!C(x,b)]$}
\DischargeRule{\Elim{\lexists}}{1}
\BinaryInfC{$\lexists[x][!C(x,b)]$}
\end{prooftree}
我们正在使用的两个假设都包含~$!B$。此时应用 \Elim{\land} 规则来分离出~$!B(a)$ 可能会很有帮助。
\begin{prooftree}
\AxiomC{$\lexists[x][(!A(x) \land !B(x)])$}
\AxiomC{$\Discharge{!A(a)
\land !B(a)}{1}$}
\RightLabel{\Elim{\land}}
\UnaryInfC{$!B(a)$}
\DeduceC{$\lexists[x][!C(x,b)]$}
\DischargeRule{\Elim{\lexists}}{1}
\BinaryInfC{$\lexists[x][!C(x,b)]$}
\end{prooftree}

我们要使用的第二个假设是~$\lforall[x][(!B(x) \lif
  !C(x,b))]$。由于它不受本征变元条件约束，我们可以用 \Elim{\lforall} 规则将 $x$ 实例化为常项符号~$a$，从而得到 $!B(a) \lif !C(a, b)$。现在我们既有 $!B(a) \lif !C(a,b)$ 又有 $!B(a)$。我们下一步应该是直接应用 \Elim{\lif} 规则。
\begin{prooftree}
\AxiomC{$\lexists[x][(!A(x) \land !B(x))]$}
\AxiomC{$\lforall[x][(!B(x) \lif !C(x,b))]$}
\RightLabel{\Elim{\lforall}}
\UnaryInfC{$!B(a) \lif !C(a,b)$}
\AxiomC{$\Discharge{!A(a)
\land !B(a)}{1}$}
\RightLabel{\Elim{\land}}
\UnaryInfC{$!B(a)$}
\RightLabel{\Elim{\lif}}
\BinaryInfC{$!C(a,b)$}
\DeduceC{$\lexists[x][!C(x,b)]$}
\DischargeRule{\Elim{\lexists}}{1}
\insertBetweenHyps{\hspace{-5em}}
\BinaryInfC{$\lexists[x][!C(x,b)]$}
\end{prooftree}
就快完成了！{} 再应用一次 \Intro{\lexists} 规则，我们就达到目标了。
\begin{prooftree}
\AxiomC{$\lexists[x][(!A(x) \land !B(x))]$}
\AxiomC{$\lforall[x][(!B(x) \lif !C(x,b))]$}
\RightLabel{\Elim{\lforall}}
\UnaryInfC{$!B(a) \lif !C(a,b)$}
\AxiomC{$\Discharge{!A(a)
\land !B(a)}{1}$}
\RightLabel{\Elim{\land}}
\UnaryInfC{$!B(a)$}
\RightLabel{\Elim{\lif}}
\BinaryInfC{$!C(a,b)$}
\RightLabel{\Intro{\lexists}}
\UnaryInfC{$\lexists[x][!C(x,b)]$}
\DischargeRule{\Elim{\lexists}}{1}
\insertBetweenHyps{\hspace{-5em}}
\BinaryInfC{$\lexists[x][!C(x,b)]$}
\end{prooftree}
由于我们在每一步都确保了本征变元条件没有被违反，因此我们可以确信这是一个正确的推导。
\end{ex}

\begin{ex}
试给出从假设 $\lforall[x][!A(x)]
\lif \lexists[y][!B(y)]$ 和 $\lnot\lexists[y][!B(y)]$ 推导公式 $\lnot\lforall[x][!A(x)]$ 的过程。像往常一样，我们将目标公式写在底部。
\begin{prooftree}
\AxiomC{}
\UnaryInfC{$\lnot\lforall[x][!A(x)]$}
\end{prooftree}
推导的最后一行是一个否定式，所以让我们尝试使用 \Intro{\lnot} 规则。这要求我们找出如何推导出一个矛盾。
\begin{prooftree}
\AxiomC{$\Discharge{\lforall[x][!A(x)]}{1}$}
\DeduceC{$\lfalse$}
\DischargeRule{\Intro{\lnot}}{1}
\UnaryInfC{$\lnot\lforall[x][!A(x)]$}
\end{prooftree}
目前进展顺利。我们可以使用 \Elim{\lforall} 规则，但这能否帮助我们达到目标尚不明确。不如先使用我们的一个假设。$\lforall[x][!A(x)] \lif \lexists[y][!B(y)]$ 与 $\lforall[x][!A(x)]$ 结合将允许我们使用 \Elim{\lif} 规则。
\begin{prooftree}
\AxiomC{$\lforall[x][!A(x)] \lif \lexists[y][!B(y)]$}
\AxiomC{$\Discharge{\lforall[x][!A(x)]}{1}$}
\RightLabel{\Elim{\lif}}
\BinaryInfC{$\lexists[y][!B(y)]$}
\DeduceC{$\lfalse$}
\DischargeRule{\Intro{\lnot}}{1}
\UnaryInfC{$\lnot\lforall[x][!A(x)]$}
\end{prooftree}
现在我们还有最后一个假设可以利用，它看起来能帮助我们通过 \Elim{\lnot} 规则得出矛盾。
\begin{prooftree}
\AxiomC{$\lnot\lexists[y][!B(y)]$}
\AxiomC{$\lforall[x][!A(x)] \lif \lexists[y][!B(y)]$}
\AxiomC{$\Discharge{\lforall[x][!A(x)]}{1}$}
\RightLabel{\Elim{\lif}}
\BinaryInfC{$\lexists[y][!B(y)]$}
\RightLabel{\Elim{\lnot}}
\BinaryInfC{$\lfalse$}
\DischargeRule{\Intro{\lnot}}{1}
\UnaryInfC{$\lnot\lforall[x][!A(x)]$}
\end{prooftree}
\end{ex}

\begin{prob}
给出推导以证明以下各式：
\begin{enumerate}
\item $\Proves (\lforall[x][!A(x)] \land \lforall[y][!B(y)]) \lif
\lforall[z][(!A(z) \land !B(z))]$.
\item $\Proves (\lexists[x][!A(x)] \lor \lexists[y][!B(y)]) \lif
\lexists[z][(!A(z) \lor !B(z))]$.
\item $\lforall[x][(!A(x) \lif !B)] \Proves \lexists[y][!A(y)] \lif !B$.
\item $\lforall[x][\lnot !A(x)] \Proves \lnot\lexists[x][!A(x)]$.
\item $\Proves \lnot\lexists[x][!A(x)] \lif \lforall[x][\lnot !A(x)]$.
\item $\Proves \lnot\lexists[x][\lforall[y][((!A(x,y) \lif \lnot
!A(y,y)) \land (\lnot !A(y,y) \lif !A(x,y)))]]$.
\end{enumerate}
\end{prob}

\begin{prob}
给出推导以证明以下各式：
\begin{enumerate}
\item $\Proves \lnot\lforall[x][!A(x)] \lif \lexists[x][\lnot!A(x)]$.
\item $(\lforall[x][!A(x)] \lif !B) \Proves \lexists[y][(!A(y) \lif !B)]$.
\item $\Proves \lexists[x][(!A(x) \lif \lforall[y][!A(y)])]$.
\end{enumerate}
（这些题目都需要用到 $\FalseCl$~规则。）
\end{prob}

\end{document}